\section{Checking comparison coverage}
\label{app:reanalysis}
We test whether the mixture gains depend on comparison coverage using the
original trained models and stored predictions. These diagnostics cover the
original 17 PDE datasets, a subset of the 21 PDE datasets in the primary
comparison. Historical evaluation outcomes had already influenced development
and scope selection, so these analyses are retrospective. The direct relative
$L^2$ fit in Appendix~\ref{app:geometry} is reported as a separate control.

\subsection{Repeating all central comparisons on fixed subsets}
\label{app:subsets}
The 14 dataset subset excludes Poisson II because of inherited numerical
defects, Cahn Hilliard I because its parameters omit realization information,
and Pressure Poisson because its coarse fields are synthetic.
The 15 dataset subset excludes Heat II and Burgers to check sensitivity to
the two B8 entries recovered by checkpoint evaluation. Their
intersection has 12 datasets. Within the revised collection, these subsets use
documented data properties and the recovered entries. None is a newly certified
benchmark. The original seven class assignment is retained, with normalization
over represented classes in each subset.

\begin{table}[h!]
\centering\small\setlength{\tabcolsep}{3pt}
\caption{\textbf{Sensitivity of the principal comparisons.} Each entry gives the
class ratio followed by the equal dataset ratio. Library and baseline denote
the respective best models identified after evaluation. Both mixtures use five fitting examples. Training budgets remain unequal.}
\label{tab:reviewsensitivity}
\resizebox{\linewidth}{!}{\begin{tabular}{lrrrr}
\toprule
Scope & $N$ & Library / baseline & \shortstack[r]{Fitted mixture /\\Selected model} & \shortstack[r]{Inverse error mixture /\\Selected model} \\
\midrule
Diagnostic set & 17 & 0.764 / 0.820 & 0.876 / 0.918 & 0.882 / 0.924 \\
Data quality subset & 14 & 0.747 / 0.810 & 0.873 / 0.900 & 0.875 / 0.908 \\
Without Heat II and Burgers & 15 & 0.790 / 0.865 & 0.872 / 0.935 & 0.872 / 0.940 \\
Intersection & 12 & 0.776 / 0.865 & 0.869 / 0.918 & 0.865 / 0.926 \\
\bottomrule
\end{tabular}
}
\end{table}

The common 15 setting comparison is especially relevant to a claim of beating
individual models. The fitted mixture has ratios 0.916 across classes and 1.020
across datasets to the hindsight best library member. The conclusion changes
with aggregation and coverage even though the fitted mixture improves on
fitting based single model selection.
Table~\ref{tab:reviewsubsets} reports all headline comparisons under each scope.

\begin{table}[p]
\centering\footnotesize
\caption{\textbf{All headline comparisons under each fixed scope.} The four
values of $N$ refer to Table~\ref{tab:reviewsensitivity}. Best library and best
baseline are retrospective minima of partition averaged errors. W, T, and L
use the one percent threshold. Methods use five fitting examples.}
\label{tab:reviewsubsets}\begin{tabular}{rlrrr}
\toprule
$N$ & Comparison & Class ratio & Dataset ratio & W / T / L \\
\midrule
17 & Best library / best additional baseline & 0.764 & 0.820 & 12/0/5 \\
17 & Fitted mixture / best library & 0.915 & 0.992 & 11/1/5 \\
17 & Inverse error mixture / best library & 0.921 & 0.999 & 11/0/6 \\
17 & Fitted mixture / best additional baseline & 0.699 & 0.813 & 10/0/7 \\
17 & Fitted mixture / Selected model & 0.876 & 0.918 & 12/0/5 \\
17 & Inverse error mixture / Selected model & 0.882 & 0.924 & 11/2/4 \\
14 & Best library / best additional baseline & 0.747 & 0.810 & 9/0/5 \\
14 & Fitted mixture / best library & 0.913 & 0.959 & 10/1/3 \\
14 & Inverse error mixture / best library & 0.915 & 0.968 & 10/0/4 \\
14 & Fitted mixture / best additional baseline & 0.682 & 0.777 & 9/0/5 \\
14 & Fitted mixture / Selected model & 0.873 & 0.900 & 11/0/3 \\
14 & Inverse error mixture / Selected model & 0.875 & 0.908 & 10/2/2 \\
15 & Best library / best additional baseline & 0.790 & 0.865 & 10/0/5 \\
15 & Fitted mixture / best library & 0.916 & 1.020 & 9/1/5 \\
15 & Inverse error mixture / best library & 0.916 & 1.026 & 10/0/5 \\
15 & Fitted mixture / best additional baseline & 0.723 & 0.882 & 8/0/7 \\
15 & Fitted mixture / Selected model & 0.872 & 0.935 & 10/0/5 \\
15 & Inverse error mixture / Selected model & 0.872 & 0.940 & 10/1/4 \\
12 & Best library / best additional baseline & 0.776 & 0.865 & 7/0/5 \\
12 & Fitted mixture / best library & 0.914 & 0.987 & 8/1/3 \\
12 & Inverse error mixture / best library & 0.909 & 0.996 & 9/0/3 \\
12 & Fitted mixture / best additional baseline & 0.709 & 0.854 & 7/0/5 \\
12 & Fitted mixture / Selected model & 0.869 & 0.918 & 9/0/3 \\
12 & Inverse error mixture / Selected model & 0.865 & 0.926 & 9/1/2 \\
\bottomrule
\end{tabular}

\end{table}

\begin{table}[p]
\centering\footnotesize
\caption{\textbf{Budget and subset sensitivity of the two mixtures.} Ratios compare with the selected model within each scope.
$N=17,14,15$ denote the historical, audit retained, and common baseline subsets.
All fits share the original models.}
\label{tab:reviewbudgets}\begin{tabular}{rrlrrr}
\toprule
$K$ & $N$ & Rule & Class ratio & Dataset ratio & W / T / L \\
\midrule
5 & 17 & Inverse error mixture & 0.882 & 0.924 & 11/2/4 \\
5 & 17 & Fitted mixture & 0.876 & 0.918 & 12/0/5 \\
5 & 14 & Inverse error mixture & 0.875 & 0.908 & 10/2/2 \\
5 & 14 & Fitted mixture & 0.873 & 0.900 & 11/0/3 \\
5 & 15 & Inverse error mixture & 0.872 & 0.940 & 10/1/4 \\
5 & 15 & Fitted mixture & 0.872 & 0.935 & 10/0/5 \\
10 & 17 & Inverse error mixture & 0.882 & 0.929 & 12/0/5 \\
10 & 17 & Fitted mixture & 0.864 & 0.900 & 12/2/3 \\
10 & 14 & Inverse error mixture & 0.885 & 0.924 & 10/0/4 \\
10 & 14 & Fitted mixture & 0.869 & 0.896 & 10/2/2 \\
10 & 15 & Inverse error mixture & 0.876 & 0.946 & 11/0/4 \\
10 & 15 & Fitted mixture & 0.862 & 0.914 & 10/2/3 \\
20 & 17 & Inverse error mixture & 0.881 & 0.923 & 12/0/5 \\
20 & 17 & Fitted mixture & 0.860 & 0.893 & 13/1/3 \\
20 & 14 & Inverse error mixture & 0.883 & 0.920 & 10/0/4 \\
20 & 14 & Fitted mixture & 0.862 & 0.888 & 11/1/2 \\
20 & 15 & Inverse error mixture & 0.875 & 0.939 & 11/0/4 \\
20 & 15 & Fitted mixture & 0.856 & 0.906 & 12/0/3 \\
\bottomrule
\end{tabular}

\end{table}

\subsection{Duplicate library entries and interpretability of weights}
The prediction source hashes identify the same complete archive for M8 and M9
on Darcy. These two entries therefore do not supply independent
prediction diversity on that dataset. The reanalysis preserves the original
library, including these entries, so it measures the stated historical procedure.
For identical models, many fitted weight allocations have the same combined
field. The inverse error mixture is not invariant to duplicating an entry, since each
copy receives weight before normalization. Effective model count measures
concentration over entries, not the number of distinct mechanisms. A deduplicated
library is a different design whose utility requires a separately specified
evaluation.

\subsection{Additional evaluation inputs for Heat I}
\label{app:additionalheat}
This supplementary experiment evaluates Heat I, one of the 22 datasets, on
additional inputs. It uses 40 new fitting pool cases and 128 new evaluation
cases with the previously trained library. Its scores are reported separately
from the main comparison tables. It evaluates the three main rules and
predates the direct relative $L^2$ fitting control.
\begin{table}[h!]
\centering\small
\caption{\textbf{Additional Heat I inputs.} Mean relative $L^2$ errors in percent. Three
fitting selections share the same 128 evaluation cases and trained models.
Selected model, inverse error mixture, and fitted mixture are the fixed rules in Section~\ref{sec:method}.}
\label{tab:fresh}\begin{tabular}{rrrrr}
\toprule
$K$ & \shortstack[r]{Selected\\model} & \shortstack[r]{Inverse\\error\\mixture} & \shortstack[r]{Fitted\\mixture} & All pairs FNO \\
\midrule
5 & 0.008638 & 0.005300 & 0.005454 & 0.011171 \\
10 & 0.008638 & 0.005286 & 0.005359 & 0.011171 \\
20 & 0.008638 & 0.005270 & 0.005296 & 0.011171 \\
\bottomrule
\end{tabular}

\end{table}
